\documentclass[11pt]{article}
\usepackage[margin=25mm]{geometry}
\usepackage{fontspec,unicode-math}
\setmainfont{Latin Modern Roman}
\setmathfont{Latin Modern Math}
\usepackage{amsmath,amssymb,mathtools}
\usepackage{xurl,hyperref}
\hypersetup{colorlinks=true,urlcolor=blue,linkcolor=blue}
\setcounter{secnumdepth}{0}
\setlength{\parindent}{0pt}
\setlength{\parskip}{5pt}
\setlength{\emergencystretch}{3em}
\newcommand{\sref}[2]{\hyperlink{source:#1:#2}{[S#2]}}
\newcommand{\cataloguescope}[1]{\paragraph{Source question.}#1}
\newcommand{\floor}[1]{\left\lfloor#1\right\rfloor}
\begin{document}
\section{Quartic Markov bases and saturated simplicial leaves}\label{statement}

\subsection{Definitions and mathematical statement}

\paragraph{Shared notation.}Standard number systems are denoted by $\mathbb N,\mathbb Z,\mathbb Q,\mathbb R,\mathbb C$; index ranges, computational models and specialized notation follow the cited problem source and the definitions below.


For an integer matrix $A$, the fiber at $b$ is $\mathcal F_b=\{x\in\mathbb Z_{\ge0}^d:Ax=b\}$. A Markov basis is a finite set of integer moves $u$ with $Au=0$ such that steps $x\mapsto x\pm u$ staying nonnegative connect every nonempty fiber. Binary hierarchical models have two levels per variable. For a move $u$, write $u=u^+-u^-$ with disjoint nonnegative parts; its degree is the total $\sum_i u_i^+$, equal to $\sum_i u_i^-$ when the total-count statistic is fixed. The task is to characterize models admitting a Markov basis with every move of degree at most four.
\paragraph{Statement recorded in the supplied source.}
Problem 28. Characterize the binary hierarchical models which have Markov bases con-sisting of moves of degree 4 or less. \sref{statement}{2}
\paragraph{Scope qualification.}
The supplied extraction includes adjacent workshop exposition; the numbered target passage is retained verbatim and that surrounding exposition is omitted.
\paragraph{Source status.}
The supplied catalogue labels this record partially solved. Its source-status metadata was checked against the supplied snapshot on 2026-10-11. This import does not independently establish current open status or mathematical correctness.
\subsection{Short English statement}
Which binary hierarchical table models can be connected using only Markov moves of degree four or less?
\subsection{Sources}
\begin{itemize}
\item[\hypertarget{source:statement:1}{S1}] Ulam AI and the UnsolvedMath Contributors, UnsolvedMath: A Curated Collection of Open Mathematics Problems, catalogue entry. Catalogue curation and metadata: CC BY 4.0. The supplied source question is retained; definitions and the English summary are editorial additions.
\url{https://huggingface.co/datasets/ulamai/UnsolvedMath}.
\item[\hypertarget{source:statement:2}{S2}] AIM Workshop Problem Lists, original problem source linked by the catalogue
\url{https://aimath.org/WWN/compalgstat/compalgstat.pdf}.
\end{itemize}
\paragraph{Preparation and review.}
Original question attributed to the cited source and UnsolvedMath contributors. Codex (GPT-6; exact serving revision not exposed) assisted with definitions, English summaries, source inspection and formatting. No new human mathematical authorship or independent proof verification is asserted. Catalogue identity checked against the complete supplied registry and existing repository statements; known flagged repetitions and obvious identical conjectures excluded. Full mathematical-equivalence/literature review has not been performed. Statement preparation and final review followed contribution prompts C and D; automated checks verify structure and identity, not open status.
\label{end:statement}
\end{document}
